\section{Progettazione Ad Alto Livello}
\label{sec:ProgettazioneAdAltoLivello}

\subsection{Front-end}
\label{sec:Front-end}
\subsubsection{HomePage}
\paragraph{Descrizione}
HomePage è il componente che guida all'utilizzatore nelle altre pagine web.
\paragraph{SottoComponenti}
HomePage è composto dai seguenti sotto-componenti:
\begin{itemize}
    \item \textbf{MenuComponent}: componente condiviso dalle pagine principali(cioè TestListPage, DatasetListPage e se stesso) permette di accedere alle pagine di TestListPage e di DatasetListPage oltre a se stesso.
    \begin{itemize}
        \item \textbf{LogoComponent}: logo del team di sviluppo.
    \end{itemize}
    
    \item \textbf{ContentComponent}: componente che descrive le informazioni riguardanti al progetto di sviluppo e al team di sviluppo.
    \item \textbf{FootterComponent}: componente che contiene le informazioni riguardanti al copyright.
\end{itemize}

\subsubsection{DataseListPage}
\label{sec:DataseListPage}
\paragraph{Descrizione}
DataSetPage è il componete che mostra le informazioni riguardanti tutti dataset salvati e temporaneo. Espone le seguenti funzionalità:
\begin{itemize}
    \item selezione/carica di un dataset salvato.
    \item modifica/rinomina di un dataset salvato.
    \item eliminazione di un dataset salvato.
    \item Duplicazione/copia di un dataset salvato.
    \item ricerca di un dataset salvato.
    \item creazione di un dataset temporaneo.
    \item Scorrimento della lista di dataset salvato.
    \item Visualizza nome dei dataset e informazioni sulla data dell'ultima modifica dei dataset.
\end{itemize}
\paragraph{SottoComponenti}
DataSetPage è composto dai seguenti sotto-componenti:
\begin{itemize}
    \item \textbf{CreaDatasetBtn}: pulsante per la creazione di un dataset temporaneo.
    \item \textbf{UploadBtn}: pulsante per caricare i file locale in formato JSON tramite una finestra, .
    \item \textbf{SearchDatasetBar}: la barra di ricerca per cercare i dataset.
    \item \textbf{DatasetListView}: mostra tutte le dataset sia salvate che temporaneo sotto forma di una lista scorrevole.
    \begin{itemize}
        \item \textbf{DatasetNameBtn}: pulsante che mostra il nome del dataset e per caricare il dataset selezionato.
        \item \textbf{DatasetChangeBtn}: pulsante per modificare il nome del dataset selezionato.
        \item \textbf{DatasetDeleteBtn}: pulsante per eliminare il dataset selezionato.
        \item \textbf{DatasetCopiaBtn}: pulsante per duplicare il dataset selezionato.
        \item \textbf{DatasetInfoText}: viene visualizzato la data della l'ultima modifica del dataset.
    \end{itemize}
\end{itemize}

\paragraph{Tracciamento dei requisiti per DataseListPage}
\begin{tabularx}{\textwidth}{|c|c|}
    \hline
    \textbf{ID} & \textbf{componente} \\
    \hline
        RFO-1.12 & CreaDatasetBtn\\
    \hline
        RFO-6.01 & DatasetNameBtn\\
    \hline
        RFO-6.02 & DatasetChangeBtn\\
    \hline
        RFO-6.03 & DatasetChangeBtn\\
    \hline
        RFO-6.04 & DatasetChangeBtn\\
    \hline
        RFO-6.04 & DatasetChangeBtn\\
    \hline
        RFO-7.01 & DatasetListView\\
    \hline
        RFO-7.02 & CreaDatasetBtn\\
    \hline
        RFO-7.03 & DatasetListView\\
    \hline
        RFO-7.04 & SearchDatasetBar\\
    \hline
        RFO-8.01 & DatasetNameBtn\\
    \hline
        RF0-8.02 & DatasetInfoText\\
    \hline
        RFO-8.03 & DatasetChangeBtn\\
    \hline
        RF0-8.04 & DatasetCopiaBtn\\
    \hline
        RF0-8.05 & DatasetDeleteBtn\\
    \hline
        RF0-8.06 & DatasetDeleteBtn\\
    \hline
        RF0-8.07 & DatasetNameBtn\\
    \hline
        RF0-8.09 & DatasetListView\\
    \hline
        RFF-13.01 & UploadBtn\\
    \hline
        RFF-13.02 & UploadBtn\\
    \hline
        RFF-13.03 & UploadBtn\\
    \hline
        RFF-13.04 & UploadBtn\\
    \hline
\end{tabularx}

\subsubsection{DataSetPageLoaded}
\label{sec:DataSetPageLoaded}
\paragraph{Descrizione}
DataSetPageLoaded è il componente che permette di visualizzare le coppie di domande-risposte del dataset selezionato tramite NameBtn.Espone le seguenti funzionalità:
\begin{itemize}
    \item Scorrimento della lista delle domande e risposte attese.
    \item Paginazione della lista delle domande e risposte attese.
    \item Modifica delle domande e risposte attese.
    \item Eliminazione di una coppia di domande e risposte attese.
    \item Ricerca delle domande e risposte tramite parole chiave.
    \item Aggiungere una coppia di domanda e risposta attesa.
    \item Eseguire un test con un LLM.
    \item Visualizza nome dei dataset e informazioni sulla data dell'ultima modifica dei dataset.
\end{itemize}
\paragraph{SottoComponenti}
DataSetPageLoaded è composto dai seguenti sotto-componenti:
\begin{itemize}
        \item \textbf{ReturnBtn}: pulsante che permette di ritornare nella pagina del DataSetPage.
        \item \textbf{DatasetNameText}: mostra nome del dataset selezionato.
        \item \textbf{SearchInfoBar}: la barra di ricerca per le domande e le risposte
        \item \textbf{SaveDatasetBtn}: il pulsante che permette di salvare/aggiornare il dataset all'interno della Database.
        \item \textbf{QAListView}: mostra tutte le coppie di domande e risposte sotto forma della lista scorrevole e della paginazione.
        \begin{itemize}
            \item \textbf{QAModificaBtn}: il pulsante che permette di modificare la coppia di domanda e risposta
            \item \textbf{QADeleteBtn}: il pulsante che permette di eliminare la coppia di domande e risposta.
        \end{itemize}

        \item \textbf{Form}: il form che permette all'utente di inserire informazioni nei campi della domanda e la risposta. 
        \begin{itemize}
            \item \textbf{CreateQABtn}: il pulsante che permette di aggiungere la coppia di domanda e risposta all'interno del dataset selezionato.
        \end{itemize}

        \item \textbf{TestBtn}:il pulsante che permette di eseguire un test con un LLM poi l'utente verrà portata nella pagina adatta.
        \begin{itemize}
            \item \textbf{QAIncompleteLink}: una pop-up che contiene tutte le domande e risposte incomplete sotto forma di un lista scorrevole coi link, viene mostrato solo in caso che ci sia presente di almeno una coppia incompleta.
            \item \textbf{ConfigLLMBtn}:una pulsante che mostra una pop-up che contiene una lista della di LLM disponibile.(parte di matteo)
        \end{itemize}

    \end{itemize}

\paragraph{Tracciamento dei requisiti per DataSetPageLoaded}
\begin{tabularx}{\textwidth}{|c|c|}
    \hline
    \textbf{ID} & \textbf{componente} \\
    \hline
        RF0-1.01 & DatasetNameText\\
    \hline
        RF0-1.02 & DatasetNameText\\
    \hline
        RF0-1.03 & QAListView\\
    \hline
        RF0-1.04 & QAListView\\
    \hline
        RF0-1.05 & QAListView\\
    \hline
        RF0-1.06 & SearchInfoBar\\
    \hline
        RF0-1.07 & CreateQABtn\\
    \hline
        RF0-1.08 & TestBtn\\
    \hline
        RF0-1.09 & QAIncompleteLink\\
    \hline
        RF0-1.10 & TestBtn\\
    \hline
        RF0-1.11 & SaveDatasetBtn\\
    \hline
        RF0-1.13 & SaveDatasetBtn\\
    \hline
        RFF-1.01 & ConfigLLMBtn\\ %parte di matteo eghosa
    \hline
        RF0-2.01 & QAListView\\
    \hline
        RF0-2.02 & QAListView\\
    \hline
        RF0-2.03 & QAListView\\
    \hline
        RF0-3.01 & QAListView\\
    \hline
        RF0-3.02 & QAListView\\
    \hline
        RF0-3.03 & QAListView\\
    \hline
        RF0-3.04 & QAModificaBtn\\
    \hline
        RF0-3.05 & QADeleteBtn\\
    \hline
        RF0-3.06 & QADeleteBtn\\
    \hline
        RF0-4.01 & QAListView\\
    \hline
        RF0-4.02 & QAListView\\
    \hline
        RF0-4.03 & QAModificaBtn\\
    \hline
        RF0-4.04 & QAModificaBtn\\
    \hline
        RF0-4.05 & QAModificaBtn\\
    \hline
        RF0-5.01 & QAIncompleteLink\\
    \hline
        RF0-5.02 & QAIncompleteLink\\
    \hline
        RF0-5.03 & QAIncompleteLink\\
    \hline
        RF0-8.08 & ReturnBtn\\
    \hline
\end{tabularx}

\subsubsection{TestListPage}
\label{sec:TestListPage}
\paragraph{Descrizione}
TestListPage è il componete che mostra le informazioni riguardanti tutti test salvati. Espone le seguenti funzionalità:
\begin{itemize}
    \item Scorrimento dei test salvati.
    \item Eliminazione dei test salvati.
    \item Ricerca dei test salvati tramite parole chiave.
    \item Caricare un test salvato.
    \item Confronto tra due test salvati.
    \item Visualizzare la data di esecuzione dei test.
\end{itemize}

\paragraph{SottoComponenti}
TestListPage è composto dai seguenti sotto-componenti:
\begin{itemize}
    \item \textbf{SearchTestBar}: la barra di ricerca per cercare i test.
    \item \textbf{TestListView}: mostra tutti i test salvati sotto forma di una lista scorrevole.
    \begin{itemize}
        \item \textbf{TestNameBtn}: pulsante che mostra il nome del test e per caricare.
        \item \textbf{TestDeletBtn}: pulsante per eliminare il test selezionato.
        \item \textbf{TestCompareBtn}: pulsante per comparare due test selezionato.
        \item \textbf{TestInfoText}: viene visualizzato la data di esecuzione del test e LLM utilizzato.
    \end{itemize}
\end{itemize}

\paragraph{Tracciamento dei requisiti per TestListPage}
\begin{tabularx}{\textwidth}{|c|c|}
    \hline
    \textbf{ID} & \textbf{componente} \\
    \hline
        RF0-9.01 & TestInfoText\\
    \hline
        RF0-9.06 & TestNameBtn\\
    \hline
        RFF-9.01 & TestNameBtn,TestInfoText\\
    \hline
\end{tabularx}

\subsubsection{TestPageLoaded}
\label{sec:TestPageLoaded}
\paragraph{Descrizione}
TestPageLoaded è il componete che mostra tutte le informazioni relativi al test selezionato ed eseguito. Espone le seguenti funzionalità:
\begin{itemize}
    \item Selezione di un test per confrontare con quello attuale.
    \item Visualizzazione del nome di test.
    \item Visualizzazione di un diagramma di dispersione.
    \item Visualizzazione dell'indice riassuntivo.
    \item Visualizzazione del diagramma di dispersione.
    \item Visualizzazione delle domande e risposte attese e ricevute.
\end{itemize}
\paragraph{SottoComponenti}
TestPageLoaded è composto dai seguenti sotto-componenti:
\begin{itemize}
    \item \textbf{NameLabel}: l'etichetta che mostra il nome del dataset.
    \item \textbf{SelectBtn}: il pulsante che mostra tutti i test salvate e che permette di selezionare un test tra di loro.
    \item \textbf{Summary}: l'indice riassuntivo.
    \item \textbf{Diagram}: visualizzazione del diagramma di dispersione.
    \item \textbf{QATestListView}: mostra tutte domande, risposte attese e risposte ricevuto dal LLM sotto forma della lista scorrevole e della paginazione.
\end{itemize}

\paragraph{Tracciamento dei requisiti per TestPageLoaded}
\begin{tabularx}{\textwidth}{|c|c|}
    \hline
    \textbf{ID} & \textbf{componente} \\
    \hline
        RFF-9.02 & SelectBtn \\
    \hline
        RFF-9.03 & QATestListView \\
    \hline
        RFF-9.04 & QATestListView \\
    \hline
        RFF-9.05 & QATestListView \\
    \hline
        RFF-9.06 & NameLabel \\
    \hline
        RFF-10.01 & Summary \\
    \hline
        RFF-10.02 & Summary \\
    \hline
        RFF-10.03 & Summary \\
    \hline
        RFF-10.04 & Summary \\
    \hline
        RFF-10.05 & Summary \\
    \hline
        RFF-10.06 & Summary \\
    \hline
        RFF-11.01 & Diagram \\
    \hline
        RFF-11.02 & QATestListView \\
    \hline
        RFF-11.03 & QATestListView \\
    \hline
        RFF-11.04 & QATestListView \\
    \hline
        RFF-12.01 & QATestListView \\
    \hline
        RFF-12.02 & QATestListView \\
    \hline
        RFF-12.03 & QATestListView \\
    \hline
        RFF-12.04 & Diagram \\
    \hline
        RFF-12.05 & QATestListView \\
    \hline
        RFF-13.01 & Diagram \\
    \hline
        RFF-13.02 & Diagram \\
    \hline
        RFF-13.03 & Diagram \\
    \hline
        RFF-13.04 & Diagram \\
    \hline
        RFF-13.05 & Diagram \\
    \hline
\end{tabularx}